\section{Numerical verification and resolution limits}
\label{sec:verification}
This section reports material conservation and solver diagnostics, then
the separate time-step and mesh comparisons, and finally two effects
that limit how precisely local profile features are resolved.

\subsection{Completion, material balance and solver diagnostics}
\label{sec:conservation}
All 15 runnable baseline cases, all 64 follow-up cases, the 19
production cases and the three full-run corrector cases completed and
passed their configured numerical acceptance checks; the scheme and
mesh runs completed with the same ledgers. This PASS
classification means only that the computation completed and passed
those checks; it is not an experimental-validation classification.
Element conservation is examined in two ways: through the solver
ledger and through an independently formed material residue. For
element $k$ the residue is
\begin{equation}
 \begin{aligned}
 R_k={}&N_k^{\rm gas}+N_k^{\rm wall}+N_k^{\rm sample}
     +N_k^{\rm net\,out}+N_k^{\rm removed}\\
     &-N_k^{\rm initial}-N_k^{\rm released},
 \end{aligned}
 \label{eq:balance}
\end{equation}
where $N_k$ are molar atom inventories and outward transport is signed.
The ledger also records the linear-solve defect and any clamping
correction separately. These terms explain bookkeeping closure, but
they are not inserted into $R_k$ when the raw material balance is
assessed.

All residues below are divided by the supplied material and taken as
the largest value over recorded steps and represented elements. The
largest raw residue is $\ProductionResidue$ for the 19 production cases
(41,472 cells, limited-linear advection) and $8.85\times10^{-13}$ for
the 60 source-driven follow-up cases; the 18,432-cell upwind baselines
had $5.97\times10^{-13}$. For the follow-up cases, the largest
accounted closure is $9.60\times10^{-16}$ and the largest recorded
solver defect is $8.64\times10^{-13}$; for the production cases they
are $1.12\times10^{-15}$ and $1.45\times10^{-12}$. The four
flow-stop cases keep their restart ledger and are excluded from this
summary of newly accumulated source material. These results support
conservation of the executed discrete model. They do not quantify
physical accuracy or discretisation error.

The baseline logs report no gas-equilibrium solver failures and a
maximum equilibrium element residue of $6.43\times10^{-16}$.
Input-validity and solver-success counters are distinct. At startup,
sign flags caused equilibrium to be skipped in up to 1,170 cells on
one step, and the maximum negative iodine diagnostic was
$1.14\times10^{-124}$ of supplied iodine. Matched five-step startup
replays identified extremely small leading-tail inventories at these
flags; no flagged cells remained at the final step. The calculations
should therefore not be described as free of negative trace values,
even though the reported amount is immaterial to the macroscopic
balance. The sum of the individual species mole-fraction maxima was at
most $9.19\times10^{-5}$. Because these maxima need not coincide, this
is an upper bound on the reactive trace loading, not a measured
mixture mole fraction at one location.

Some paired wall-coupling iterations did not meet the requested
$10^{-16}$ outer criterion. This is reported separately from
equilibrium failure, the final transport defect and the inventory-bound
checks. To test its effect, matched restarts of four representative
cases at 59 s compared three and eight correctors, and outer
tolerances of $10^{-16}$ and $10^{-14}$, over 20 time steps (40 ms).
The largest species-field change, scaled by its control maximum, was
below $3.6\times10^{-8}$, and the changes in new wall increments were
negligible. Three cases were then repeated over the full 60 s with
eight instead of three correctors (18,432 cells, upwind). The iodine
profiles of pure PbI$_2$ and of the steel Q2 case at 0.1\% saturation
change by $\IterPure\%$ and $\IterSteelLow\%$, that is, at round-off;
the steel case at 1\% saturation changes by \IterSteelHigh\%. With
eight correctors the largest outer initial residual is still
$2.4\times10^{-15}$--$4.4\times10^{-15}$, and nearly every step ends
at the corrector limit, as with three. The requested $10^{-16}$ lies
below the attainable residual floor, and the iteration error of the
reported profiles is negligible next to the time-step and spatial
errors below. No manufactured-solution test is presented.
\Cref{sec:lobresco} compares the model with published measurements
and published model results; it is not a code-to-code verification
with identical inputs.

\subsection{Separate time-step and mesh comparisons}
For profiles on the common 1 cm bins, let $a_b$ be the deposited
iodine atoms in bin $b$ for a comparison case and $a_b^0$ the
corresponding baseline inventory. We report the profile change
$E_{L1}$ and the total change $\Delta N_I$:
\begin{equation}
 E_{L1}=100\frac{\sum_b|a_b-a_b^0|}{\sum_b a_b^0},
 \qquad
 \Delta N_I=100\frac{\sum_b a_b-\sum_b a_b^0}{\sum_b a_b^0}.
 \label{eq:numeric-l1}
\end{equation}
Both measures retain changes in amount. The figures also show
normalised shape curves; these must not be confused with the
unnormalised profile change. Source, physical coefficients,
equilibrium engine and executable are fixed within each numerical
comparison.

\Cref{tab:numerical} separates time-step from mesh and scheme effects.
For successive mesh levels the reference $a_b^0$ is the finer level.
Reducing the time step from 2 to 1 ms on 18,432 cells changes the
iodine profile by \TimeMin--0.3228\%. With the upwind scheme of the
original campaign, the 18,432 to 41,472-cell differences of the four
follow-up cases are \MeshMin--\MeshMax\%, and the shifted-BiI steel
case changes by essentially the same amount on the first two
refinements (\SteelUpOne\% and \SteelUpTwo\%). Total deposited iodine
changes by less than 0.014\% from 18,432 to 41,472 cells, and the
centroid moves by at most 0.0652 cm. Integral quantities are thus
stable while local profile shapes still differ appreciably.
\begin{table*}[tb]
\centering
\caption{Iodine profile changes $E_{L1}$ (percent) between successive
mesh levels at 2 ms, relative to the finer level, for three advection
schemes, and the effect of halving the time step on 18,432 cells with
upwind advection. All steel cases are Q2; the shifted-BiI case uses
1\% metal saturation. The upwind levels of the two lower steel cases
are the follow-up cases of \cref{tab:study}. A dash marks a combination
that was not run.}
\label{tab:numerical}
\small\begin{tabular}{llrrrr}
\toprule
Case & Scheme & 8,192$\to$18,432 & 18,432$\to$41,472 & 41,472$\to$93,312 & 2$\to$1 ms \\
\midrule
Pure PbI$_2$ & upwind & 6.57 & 5.06 & 3.68 & 0.0043 \\
 & limitedLinear & 5.70 & 2.15 & -- & -- \\
 & van Leer & 4.94 & 2.23 & -- & -- \\
\midrule
Steel, shifted BiI & upwind & 10.47 & 10.58 & 2.13 & 0.1536 \\
 & limitedLinear & 15.35 & 8.95 & 1.98 & -- \\
 & van Leer & 16.83 & 9.17 & -- & -- \\
\midrule
Steel Q2, 0.1\% & upwind & 9.02 & 5.26 & 7.30 & 0.2326 \\
\midrule
Steel Q2, 1\% & upwind & 9.72 & 2.06 & 6.68 & 0.3228 \\
\bottomrule
\end{tabular}

\end{table*}

\begin{figure*}[tb]
\centering
\includegraphics{refinement.pdf}
\caption{Mesh sequences of the iodine deposition profile for pure
PbI$_2$ (top) and the shifted-BiI steel case at 1\% metal saturation
(bottom), with first-order upwind (left) and limited-linear (right)
advection. Each curve is normalised over the common scanned interval
so that shapes can be compared; the unnormalised differences are given
in \cref{tab:numerical}. Grey histograms are the reconstructed
experimental iodine scans \cite{Liu2025chemical} and include any KI
contribution.}
\label{fig:refinement}
\end{figure*}

The scheme and mesh study (\cref{tab:numerical,fig:refinement})
addresses the shifted-BiI steel case directly. With the limited-linear
scheme its iodine profile changes by \SteelLLOne, \SteelLLTwo{} and
\SteelLLThree\% over the three refinements to 93,312 cells; with
upwind, by \SteelUpOne, \SteelUpTwo{} and \SteelUpThree\%. Pure
PbI$_2$ changes by \PureLLOne{} and \PureLLTwo\% with limited-linear,
against \PureUpOne{} and \PureUpTwo\% with upwind. The van Leer scheme
gives successive changes close to those of limited-linear and differs
from it by at most \VLSteelMax\% (steel) and \VLPureMax\% (pure
PbI$_2$) at equal mesh, and by \VLSteelThree\% and \VLPureThree\% on
41,472 cells. Between 41,472 and 93,312 cells, the 3 cm smoothed
profile of the steel case changes by \SteelSmoothThree\% and its total
deposit by \SteelTotalThree\%; over all four levels its centroid stays
within \SteelCentroidRange{} mm and the wall temperature at the PbI$_2$
centroid within \SteelCentroidTRange{} K.

Three-grid scalar checks follow the distinction between refinement
trends and an established asymptotic range
\cite{NASAGridConvergence,Roache1994perspective,Celik2008procedure}.
For the earlier upwind follow-up triplets, 17 of the 34 nontrivial or
zero-valued combinations of species, total, centroid and width are
oscillatory, four have nondecreasing differences, one is unresolved at
roundoff and 12 admit only conditional Richardson/GCI estimates. For
the limited-linear steel sequence, the iodine total, centroid and
width oscillate on both triplets, and the successive profile changes
give apparent orders of \SteelOrderOne{} and \SteelOrderTwo. These
orders are inconsistent, so no formal grid convergence index is
claimed. We quote the measured 41,472 to 93,312-cell change of
\SteelLLThree\% as the spatial uncertainty indicator of the production
profiles. If the last change were extrapolated at the last apparent
order, the 41,472-cell profile would lie about \SteelErrApparent\% from
the mesh limit; assuming second order, about \SteelErrSecond\%. These
are indicative estimates for one case, not confirmed uncertainty
bounds, and they do not establish mesh independence of local profile
shapes in other cases.

\subsection{Binned maxima and numerical diffusion}
A deposition maximum picked on a 1 cm grid is fragile when adjacent
bins have similar heights. In pure PbI$_2$ the 18,432-cell upwind
maximum is at 37 cm and 350.7\,$^\circ$C, while the 36 cm bin is only
about 0.05\% lower. On the 41,472-cell upwind mesh the maximum moves to
36 cm and 366.5\,$^\circ$C, although the centroid shifts by only
$-0.0585$ cm. With the production scheme (41,472 cells, limited-linear)
the binned maximum stays at 36 cm and 366.4--366.5\,$^\circ$C on all
three levels, and on 41,472 cells the 37 cm bin is 23\% lower, yet the maximum of the 3 cm
smoothed profile is at 37 cm and 350.5\,$^\circ$C, and the wall
temperature at the deposit centroid (37.11 cm) is 348.7\,$^\circ$C.
In the limited-linear steel sequence the PbI$_2$ maximum alternates
between two nearly equal bins, 50 and 51 cm (306.3, 288.1, 306.7 and
288.1\,$^\circ$C on the four levels), while the 3 cm smoothed maximum
stays at 51 cm and 288.0--288.1\,$^\circ$C and the centroid temperature
at 290.6--291.5\,$^\circ$C. In the upwind steel Q2 0.1\% case, the
dominant BiI$_3$ bin moves from 57 to 58 cm, and its reported
temperature changes from about 166.4 to 138.8\,$^\circ$C. A one-bin
switch can therefore change the peak temperature by nearly 28 K without
a comparable displacement of the profile as a whole. Peak temperature
alone is an insufficient convergence metric; \cref{tab:peaks} therefore
also reports the smoothed and centroid values.

First-order upwind advection, used in the original campaign and the
follow-up cases, smooths the profiles. The local one-dimensional
estimate $D_{\rm num}\simeq|u_x|\Delta x/2$ is comparable to or greater
than the prescribed $10^{-4}\units{m^2\,s^{-1}}$ diffusivity in parts of
the carriers. This is a Taylor-series screening estimate, not a
measured truncation error for the radial, variable-density coupled
problem, but the scheme study confirms its consequence. At equal mesh,
upwind differs from limited-linear by \UpLLOne, \UpLLTwo, \UpLLThree{}
and \UpLLFour\% for the steel case on the four levels and by
\PureUpLLMin--\PureUpLLMax\% for pure PbI$_2$, much more than the two
bounded schemes differ from each other. Upwind numerical diffusion is
therefore the dominant spatial error of the original campaign, and it
is why halving the time step could not establish convergence of
profile widths. Across the 19 production cases, the production profiles
differ from the 18,432-cell upwind profiles by 9.1--22.3\%. The
follow-up sensitivities of \cref{sec:sensitivity} remain upwind
differences on 18,432 cells, taken against baselines with the same
numerical error.

\subsection{Dependence on comparison-bin width}
\label{sec:bin-audit}
To distinguish sensitivity of the reported observable from a change in
the CFD mesh, the saved 1 cm iodine inventories of the original upwind campaign
were summed into groups
of two and four consecutive bins. Every possible integer-bin offset was
tested for each width, retaining partial end groups and the full deposited
inventory. No CFD field was recomputed and no sub-bin reconstruction was
introduced. \Cref{tab:bin-audit} gives the fine/medium inventory differences
for these coarser observables. Ranges reflect bin origin, not statistical
uncertainty.
\begin{table*}[tbp]
\centering
\caption{Original upwind fine/medium iodine inventory-profile differences
(percent, 41,472 versus 18,432 cells) after
aggregation of native bins. The 2 and 4 cm ranges cover all respective
integer-bin offsets. $W_1$ uses the original 1 cm bin centres and separately
normalised profiles; it measures redistribution distance, not position error
against the experiment.}
\label{tab:bin-audit}
\small\begin{tabular}{lrrrr}
\toprule
Case & 1 cm & 2 cm & 4 cm & $W_1$ (cm) \\
\midrule
Pure PbI2 & 5.06 & 4.09--4.81 & 1.22--4.81 & 0.0610 \\
Steel Q2, 0.1\% & 5.26 & 1.77--4.12 & 0.51--2.45 & 0.0311 \\
Steel Q2, 1\% & 2.06 & 0.85--1.43 & 0.11--1.31 & 0.0115 \\
Steel Q2, shifted BiI & 10.58 & 3.37--10.02 & 0.52--10.02 & 0.0711 \\
\bottomrule
\end{tabular}

\end{table*}

For the shifted-BiI case, a 4 cm comparison gives differences from
0.52\% to 10.02\% depending on the grouping origin. Thus a favourable
aggregation can hide a substantial local difference. Summing signed
bin differences before taking their absolute values necessarily allows
cancellation; reduced aggregated L1 is not additional convergence evidence.
The largest difference remains appreciable even at 4 cm resolution.

As a complementary position-sensitive measure, let $p_b$ and $q_b$ be
the fine and medium iodine inventories normalised separately over the
whole column, and $x_b$ their common native bin centres. The
one-dimensional transport distance for these discrete distributions is
\begin{equation}
 W_1=\sum_{b=1}^{B-1}\left|\sum_{j=1}^{b}(p_j-q_j)\right|
                  (x_{b+1}-x_b).
 \label{eq:binned-wasserstein}
\end{equation}
It is the minimum mean distance needed to redistribute one normalised
binned inventory into the other. Values of 0.0115--0.0711 cm indicate
small average spatial redistribution between the two tested meshes,
despite larger local L1 differences. This is a distance between computed
binned profiles, not a bound on continuum error or evidence that the
experimental position uncertainty has been resolved.
